\documentclass[conference]{IEEEtran}
\IEEEoverridecommandlockouts

\usepackage{cite}
\usepackage{amsmath,amssymb,amsfonts}
\usepackage{graphicx}
\usepackage{textcomp}
\usepackage{xcolor}
\usepackage{booktabs}
\usepackage{hyperref}

\begin{document}

\title{
A Multi-Signal Probabilistic Temporal Risk Accumulation Framework\\
for AI-Based Online Proctoring Under Strict One-Shot Identity Constraints
}

\author{
\IEEEauthorblockN{Shyam Saran P}
\IEEEauthorblockA{
Department of Information Technology\\
SSN College of Engineering\\
Chennai, India\\
shyamsaran348@gmail.com}
}

\maketitle

% ============================================================
\begin{abstract}
We present a multi-signal probabilistic temporal risk accumulation framework for AI-based online proctoring under strict one-shot identity constraints. The system integrates an Identity Extractor (IDE), Identity Instability Model (IIM), Presence and Attentiveness Model (PAM), Long-term Drift Detector (LDD), and a GRU-based Risk Fusion Engine (RFE). All signals are continuous and temporally modeled without heuristic thresholds or rule-based overrides. Evaluation on procedurally generated temporally realistic sessions demonstrates strong discrimination (ROC-AUC = 0.9987), effective calibration (ECE = 0.0072), and smooth risk trajectories. Results confirm that long-horizon drift modeling contributes uniquely to detection of gradual impersonation scenarios.
\end{abstract}

\begin{IEEEkeywords}
Online Proctoring, Temporal Modeling, Risk Fusion, One-Shot Identity, Deep Learning
\end{IEEEkeywords}

% ============================================================
\section{Introduction}

Online examination integrity is increasingly challenged by remote assessment environments. Traditional rule-based systems rely on threshold-triggered alerts, while shallow classifiers collapse temporal behavior into binary decisions. We instead formulate online examination integrity as a multi-signal temporal probabilistic inference problem under strict one-shot identity constraints.

The proposed framework operates in three conceptual layers:
\begin{enumerate}
\item Frame-level feature extraction (IDE, PAM)
\item Temporal modeling (IIM, LDD)
\item Session-level probabilistic risk accumulation (RFE)
\end{enumerate}

% ============================================================
\section{System Overview and Workflow}

\begin{figure}[t]
\centering
\includegraphics[width=\linewidth]{fig7_workflow.png}
\caption{Operational workflow from enrollment to session-level risk estimation.}
\label{fig:workflow}
\end{figure}

Enrollment produces a single immutable embedding $e_0$. During monitoring, each frame $f_t$ produces probe embedding $e_t$, similarity $S_t$, instability $I_t$, presence $P_t$, and drift $D_t$. The Risk Fusion Engine (RFE) updates risk $\rho_t$ over time, with final session risk $\rho_T$ committed without threshold overrides.

% ============================================================
\section{Mathematical Formulation}

Given enrollment embedding $e_0$ and probe embedding $e_t$:

\begin{equation}
S_t = e_t \cdot e_0
\end{equation}

Delta embedding:
\begin{equation}
\delta_t = e_t - e_0
\end{equation}

Risk vector:
\begin{equation}
r_t = [S_t, I_t, P_t, D_t]
\end{equation}

RFE GRU update:
\begin{equation}
\rho_t = \sigma(W_h h_t + b)
\end{equation}

Training uses session-level binary cross-entropy on $\rho_T$ only.

% ============================================================
\section{System Architecture}

\begin{figure*}[t]
\centering
\includegraphics[width=\textwidth]{fig6_system_architecture.png}
\caption{Multi-signal architecture under strict one-shot identity constraint.}
\label{fig:architecture}
\end{figure*}

The Identity Extractor (IDE) computes embeddings. The Identity Instability Model (IIM) detects short-term fluctuations in similarity. The Presence and Attentiveness Model (PAM) models behavioral consistency. The Long-term Drift Detector (LDD) captures cumulative embedding displacement. The Risk Fusion Engine (RFE) integrates all signals temporally.

% ============================================================
\section{Experimental Setup}

Evaluation uses 5,000 procedurally generated sessions (120 frames each) with AR(1) temporal correlation and overlapping Gaussian distributions. Models are trained using Adam optimizer with cosine learning rate schedule.

% ============================================================
\section{Results}

\begin{figure}[t]
\centering
\includegraphics[width=\linewidth]{fig1_ablation.png}
\caption{Signal ablation analysis (ROC-AUC).}
\label{fig:ablation}
\end{figure}

\begin{figure}[t]
\centering
\includegraphics[width=\linewidth]{fig2_baselines.png}
\caption{Baseline comparison (ROC-AUC and Brier score).}
\label{fig:baseline}
\end{figure}

\begin{figure}[t]
\centering
\includegraphics[width=\linewidth]{fig3_trajectories.png}
\caption{Risk trajectories across session archetypes.}
\label{fig:trajectories}
\end{figure}

\begin{figure}[t]
\centering
\includegraphics[width=\linewidth]{fig4_calibration.png}
\caption{Calibration reliability comparison.}
\label{fig:calibration}
\end{figure}

The full four-signal configuration achieves ROC-AUC 0.9987. Removal of LDD produces the largest performance degradation, indicating unique contribution of long-horizon modeling.

% ============================================================
\section{Discussion}

Temporal modeling enables accumulation of weak but consistent deviations. The GRU-based RFE achieves superior calibration and smoother trajectories compared to non-temporal baselines.

% ============================================================
\section{Ethical Considerations}

The system outputs calibrated probabilities rather than binary decisions. High-risk sessions should trigger human review rather than automatic disqualification. The one-shot constraint reduces persistent biometric storage.

% ============================================================
\section{Limitations}

Evaluation relies on procedurally generated sessions. Real-world annotated datasets are required for further validation. ResNet-50 latency may require backbone substitution for real-time CPU deployment.

% ============================================================
\section{Conclusion}

We presented a probabilistic, temporally modeled proctoring framework integrating identity, behavioral, and drift signals under strict one-shot constraints. Experimental results demonstrate strong discrimination and calibration while preserving architectural invariants.

% ============================================================
\begin{thebibliography}{00}

\bibitem{facenet} F. Schroff, D. Kalenichenko, and J. Philbin, ``FaceNet: A Unified Embedding for Face Recognition,'' in CVPR, 2015.

\bibitem{arcface} J. Deng et al., ``ArcFace: Additive Angular Margin Loss for Deep Face Recognition,'' in CVPR, 2019.

\bibitem{cosface} H. Wang et al., ``CosFace: Large Margin Cosine Loss for Deep Face Recognition,'' in CVPR, 2018.

\end{thebibliography}

\end{document}
